\settopicname{Nonlinear Functions}
\chapter{Nonlinear Functions}
\topicopen{Advanced Math}{Analyze, interpret and produce nonlinear functions
in one variable --- quadratic, exponential and polynomial --- and identify
what an equivalent form or a transformation reveals about the graph.}

\begin{essentials}
Quadratics come in three forms, and each one hands you something different.
Standard $f(x)=ax^2+bx+c$: $y$-intercept $(0,c)$, axis of symmetry
$x=-\frac{b}{2a}$. Factored $f(x)=a(x-r)(x-s)$: $x$-intercepts $r$ and $s$,
axis of symmetry halfway between them at $x=\frac{r+s}{2}$. Vertex
$f(x)=a(x-h)^2+k$: vertex $(h,k)$, a minimum when $a>0$ and a maximum when
$a<0$. The graph meets the $x$-axis twice when $b^2-4ac>0$ and exactly once
when $b^2-4ac=0$.

Exponential $f(x)=ab^{\,x}$: $a=f(0)$ is the initial amount and $b$ is the
factor per unit of $x$ --- growth when $b>1$, decay when $0<b<1$ --- with a
change of $(b-1)\cdot 100\%$ per unit. In $f(x)=ab^{\,x/n}$ the factor $b$
applies once every $n$ units, so $b=2$ makes $n$ the doubling time.

A polynomial's zeros are the values that make its factors zero, and its end
behaviour follows the leading term. Transformations: $f(x)+k$ shifts up $k$,
$f(x+k)$ shifts left $k$, and $-f(x)$ reflects across the $x$-axis.
\end{essentials}

\modulehead{Module 1}{Questions 1--22}
\begin{multicols*}{2}

\question{The function $f$ is defined by $f(x)=(x-3)(x+5)$. Which of the
following is a zero of $f$?
\choices{\item $-5$ \item $-3$ \item $5$ \item $8$}}

\qfr{The function $f$ is defined by $f(x)=2x^2-7x+5$. The graph of $y=f(x)$
in the $xy$-plane has its $y$-intercept at $(0,k)$. What is the value of $k$?}

\question{The function $f$ is defined by $f(x)=(x-4)^2+9$. What is the
minimum value of $f$?
\choices{\item $-9$ \item $4$ \item $9$ \item $13$}}

\question{The function $f$ is defined by $f(x)=250(1.08)^x$. What is the
value of $f(0)$?
\choices{\item $1.08$ \item $250$ \item $258$ \item $270$}}

\qfr{The function $f$ is defined by $f(x)=(x+2)(x-8)$. The axis of symmetry
of the graph of $y=f(x)$ in the $xy$-plane is the line $x=c$. What is the
value of $c$?}

\question{In the $xy$-plane, the graph of $y=f(x)+6$ is the graph of
$y=f(x)$ shifted in which of the following ways?
\choices{\item 6 units up \item 6 units down \item 6 units left
\item 6 units right}}

\question{The function $f$ is defined by $f(x)=6(2)^x$. What is the value of
$f(3)$?
\choices{\item $18$ \item $36$ \item $48$ \item $64$}}

\qfr{The polynomial function $p$ is defined by $p(x)=x(x-4)(x+7)$. The graph
of $y=p(x)$ in the $xy$-plane crosses the $x$-axis at three points. What is
the greatest of the three $x$-coordinates?}

\question{Which of the following is equivalent to
\[ x^2-6x+11\,? \]
\choices{\item $(x-3)^2+2$ \item $(x-3)^2+11$ \item $(x+3)^2+2$
\item $(x-6)^2+2$}}

\question{The number of insects in a colony is modeled by
$P(t)=1200(1.35)^t$, where $t$ is the number of years since the study began.
What is the best interpretation of $1200$ in this context?
\choices{\item The number of insects when the study began
\item The number of insects added each year
\item The number of years the study lasted
\item The percent increase in insects each year}}

\qfr{The function $f$ is defined by $f(x)=2x^2-12x+7$. What is the
$x$-coordinate of the vertex of the graph of $y=f(x)$ in the $xy$-plane?}

\question{The value, in dollars, of an investment after $t$ years is modeled
by $f(t)=5000(1.04)^t$. By what percent does the value of the investment
increase each year?
\choices{\item $0.04\%$ \item $4\%$ \item $40\%$ \item $104\%$}}

\question{The function $g$ is defined by $g(x)=f(x+3)$. In the $xy$-plane,
the graph of $y=g(x)$ is the graph of $y=f(x)$ shifted in which of the
following ways?
\choices{\item 3 units left \item 3 units right \item 3 units up
\item 3 units down}}

\qfr{What is the positive solution to the equation $x^2+3x-40=0$?}

\question{The polynomial function $p$ is defined by $p(x)=-2x^3+5x-1$. As
$x$ increases without bound, which of the following describes the behaviour
of $p(x)$?
\choices{\item $p(x)$ increases without bound
\item $p(x)$ decreases without bound
\item $p(x)$ approaches $0$
\item $p(x)$ approaches $-1$}}

\question{The function $f$ is defined by $f(x)=x^2-2x-24$. Which of the
following displays the $x$-intercepts of the graph of $y=f(x)$ as constants
or coefficients?
\choices{\item $(x-6)(x+4)$ \item $(x+6)(x-4)$ \item $(x-4)(x-6)$
\item $(x+4)(x+6)$}}

\qfr{The function $f$ is defined by $f(t)=3(2)^{t/5}$. For what value of $t$
does $f(t)=24$?}

\question{In the given equation, $k$ is a positive constant.
\[ y=x^2+kx+36 \]
The graph of the equation in the $xy$-plane has exactly one $x$-intercept.
What is the value of $k$?
\choices{\item $6$ \item $12$ \item $18$ \item $36$}}

\question{The function $f$ is defined by $f(x)=(x-7)(x+1)$. The graph of
$y=f(x)$ in the $xy$-plane has its minimum at the point $(a,b)$. What is the
value of $b$?
\choices{\item $-16$ \item $-12$ \item $-7$ \item $3$}}

\question{The mass, in grams, of a sample of a substance $t$ days after it
was first measured is modeled by
\[ g(t)=64\left(\tfrac{1}{2}\right)^{t/18}. \]
What is the best interpretation of $18$ in this context?
\choices{\item The mass of the sample is cut in half every 18 days
\item The sample loses 18 grams each day
\item The sample is completely gone after 18 days
\item The mass of the sample was 18 grams when it was first measured}}

\question{The function $f$ is defined by $f(x)=2x^2+bx+3$, where $b$ is a
constant. The axis of symmetry of the graph of $y=f(x)$ in the $xy$-plane is
the line $x=-3$. What is the value of $b$?
\choices{\item $-12$ \item $-6$ \item $6$ \item $12$}}

\qfr{The function $h$ is defined by $h(x)=(x-a)(x-6)$, where $a$ is a
constant. The axis of symmetry of the graph of $y=h(x)$ in the $xy$-plane is
the line $x=1$. What is the value of $a$?}

\end{multicols*}

\modulehead{Module 2}{Questions 1--22}
\begin{multicols*}{2}

\question{Which of the following is equivalent to
\[ x^2+8x+3\,? \]
\choices{\item $(x+4)^2-13$ \item $(x+4)^2+3$ \item $(x-4)^2-13$
\item $(x+8)^2-13$}}

\qfr{The balance, in dollars, of an account $t$ years after it was opened is
modeled by $B(t)=400(1.05)^t$. What is the balance, in dollars, of the
account 2 years after it was opened?}

\question{The number of bacteria in a culture is modeled by $n(h)=500(3)^h$,
where $h$ is the number of hours since the culture was prepared. What is the
best interpretation of $3$ in this context?
\choices{\item The number of bacteria triples each hour
\item The number of bacteria increases by 3 each hour
\item There were 3 bacteria when the culture was prepared
\item The number of bacteria doubles every 3 hours}}

\question{The polynomial function $p$ is defined by $p(x)=x^3-9x$. What is
the greatest zero of $p$?
\choices{\item $-3$ \item $0$ \item $3$ \item $9$}}

\qfr{The function $f$ is defined by $f(x)=x^2-3x$, and the function $g$ is
defined by $g(x)=f(x)+7$. What is the value of $g(4)$?}

\question{In the $xy$-plane, the graph of $y=-f(x)$ is the graph of
$y=f(x)$ reflected across which of the following?
\choices{\item The $x$-axis \item The $y$-axis \item The line $y=x$
\item The origin}}

\question{The value, in dollars, of a machine $t$ years after it was
purchased is modeled by $v(t)=18000(0.88)^t$. What is the best
interpretation of $0.88$ in this context?
\choices{\item Each year the value of the machine is 88\% of its value the
year before
\item The machine loses 88 dollars in value each year
\item The value of the machine decreases by 88\% each year
\item The machine will be worth nothing after 88 years}}

\qfr{The function $f$ is defined by $f(x)=x^2-14x+k$, where $k$ is a
constant. The minimum value of $f$ is $5$. What is the value of $k$?}

\question{Which of the following is equivalent to
\[ (2x-5)(x+3)\,? \]
\choices{\item $2x^2+x-15$ \item $2x^2-x-15$ \item $2x^2+11x-15$
\item $2x^2+x+15$}}

\question{The polynomial function $q$ is defined by $q(x)=3x^4-7x^2+2$.
Which of the following describes the end behaviour of the graph of $y=q(x)$
in the $xy$-plane?
\choices{\item $q(x)$ increases without bound as $x$ increases and as $x$
decreases
\item $q(x)$ decreases without bound as $x$ increases and as $x$ decreases
\item $q(x)$ increases without bound as $x$ increases and decreases without
bound as $x$ decreases
\item $q(x)$ approaches 2 as $x$ increases and as $x$ decreases}}

\qfr{The population of a town is modeled by $P(t)=8000(2)^{t/12}$, where $t$
is the number of years after the year 2000. For what value of $t$ does the
model predict a population of $32{,}000$?}

\question{The function $f$ is defined by $f(x)=x^2+bx-15$, where $b$ is a
constant. The graph of $y=f(x)$ in the $xy$-plane has $x$-intercepts at
$(-5,0)$ and $(3,0)$. What is the value of $b$?
\choices{\item $-8$ \item $-2$ \item $2$ \item $8$}}

\question{The function $f$ is defined by $f(x)=x^2+1$, and the function $g$
is defined by $g(x)=f(x-4)$. What is the value of $g(6)$?
\choices{\item $5$ \item $17$ \item $37$ \item $41$}}

\qfr{The height, in feet, of a ball $t$ seconds after it is thrown upward is
modeled by $h(t)=-16t^2+64t$. What is the maximum height, in feet, reached
by the ball?}

\question{The function $f$ is defined by $f(x)=a(x-2)(x-6)$, where $a$ is a
positive constant. Which statement about the graph of $y=f(x)$ in the
$xy$-plane must be true?
\choices{\item The graph reaches its minimum value when $x$ equals 4
\item The graph reaches its maximum value when $x$ equals 4
\item The graph reaches its minimum value when $x$ equals 2
\item The graph never intersects the $x$-axis}}

\question{The function $f$ is defined by $f(x)=2x^2-16x+k$, where $k$ is a
constant. The minimum value of $f$ is $-6$. What is the value of $k$?
\choices{\item $10$ \item $22$ \item $26$ \item $38$}}

\qfr{The function $g$ is defined by $g(x)=a(x+3)(x-7)$, where $a$ is a
constant. The graph of $y=g(x)$ in the $xy$-plane passes through the point
$(1,-48)$. What is the value of $a$?}

\question{The number of members of a club $t$ months after it opened is
modeled by $f(t)=120(1.5)^{t/4}$. What is the best interpretation of $1.5$
in this context?
\choices{\item The number of members increases by a factor of 1.5 every 4
months
\item The number of members increases by a factor of 1.5 each month
\item The club gains 1.5 members each month
\item The club had 1.5 members when it opened}}

\question{The polynomial function $p$ is defined by $p(x)=(x+2)^2(x-5)$.
Which statement about the graph of $y=p(x)$ in the $xy$-plane must be true?
\choices{\item The graph touches the $x$-axis at $x=-2$ and crosses it at
$x=5$
\item The graph crosses the $x$-axis at three distinct points
\item The graph has a $y$-intercept at $(0,20)$
\item The graph decreases without bound as $x$ increases without bound}}

\qfr{The function $h$ is defined by $h(x)=x^2-8x+3$, and the function $k$ is
defined by $k(x)=h(x+2)$. What is the $x$-coordinate of the vertex of the
graph of $y=k(x)$ in the $xy$-plane?}

\question{The function $f$ is defined by $f(x)=3(4)^x$. Which of the
following is equivalent to
\[ f(x+1)\,? \]
\choices{\item $12(4)^x$ \item $3(4)^x+1$ \item $7(4)^x$
\item $3(4)^x+4$}}

\question{The function $f$ is defined by $f(x)=x^2-6x+13$, and $k$ is a
constant. The graph of $y=f(x)+k$ in the $xy$-plane has exactly one
$x$-intercept. What is the value of $k$?
\choices{\item $-13$ \item $-4$ \item $4$ \item $13$}}

\end{multicols*}
